\section{Математические модели в механике деформируемого твёрдого тела}
\label{sec:solid-mechanics-models}

\subsection{Основные переменные и кинематика деформирования}

Механика деформируемого твёрдого тела описывает движение и напряжённо--
деформированное состояние сплошного тела. В лагранжевом описании материальная
точка задаётся радиус-вектором $\mathbf R$ в исходной конфигурации, а её положение
в текущий момент времени имеет вид
\[
    \mathbf r=\mathbf r(\mathbf R,t)
    =\mathbf R+\mathbf u(\mathbf R,t),
\]
где $\mathbf u$ --- вектор перемещений. Градиент деформации
\[
    \mathbf F=\frac{\partial \mathbf r}{\partial \mathbf R}
    =\mathbf I+\nabla_{\!R}\mathbf u
\]
характеризует локальное изменение длин и углов.

При конечных деформациях удобно использовать тензор деформаций Грина
\[
    \mathbf E=\frac12\left(\mathbf F^{T}\mathbf F-\mathbf I\right).
\]
В линейной теории предполагаются малые перемещения и малые градиенты
перемещений. Тогда квадратичными членами можно пренебречь и тензор малых
деформаций принимает вид
\[
    \boldsymbol\varepsilon
    =\frac12\left(\nabla\mathbf u+(\nabla\mathbf u)^T\right),
    \qquad
    \varepsilon_{ij}=\frac12
    \left(\frac{\partial u_i}{\partial x_j}
          +\frac{\partial u_j}{\partial x_i}\right).
\]
Диагональные компоненты описывают относительные удлинения, а недиагональные
--- сдвиговые деформации.

\subsection{Тензор напряжений и уравнения баланса}

Внутренние силы в сплошной среде описываются тензором напряжений Коши
$\boldsymbol\sigma$. Вектор напряжений, действующий на площадку с единичной
нормалью $\mathbf n$, определяется формулой Коши
\[
    \mathbf t_{\mathbf n}=\boldsymbol\sigma\mathbf n.
\]
Из закона сохранения момента количества движения в классической среде без
моментных напряжений следует симметрия
\[
    \boldsymbol\sigma=\boldsymbol\sigma^T.
\]

Основное динамическое уравнение --- локальный закон баланса количества
движения:
\[
    \rho\,\ddot{\mathbf u}
    =\operatorname{div}\boldsymbol\sigma+\rho\mathbf f,
\]
где $\rho$ --- плотность, $\mathbf f$ --- объёмная сила, отнесённая к единице
массы. В квазистатическом приближении инерционным членом пренебрегают:
\[
    \operatorname{div}\boldsymbol\sigma+\rho\mathbf f=0.
\]
Таким образом, математическая модель состоит как минимум из кинематических
соотношений, уравнений баланса и определяющего закона материала.

\subsection{Определяющие соотношения. Линейная теория упругости}

Определяющее соотношение связывает напряжения с деформациями и тем самым
задаёт модель материала. Для линейно-упругого тела общий закон Гука записывается
в тензорной форме
\[
    \sigma_{ij}=C_{ijkl}\varepsilon_{kl},
\]
где $C_{ijkl}$ --- тензор упругих модулей.

Для однородного изотропного материала закон Гука имеет вид
\[
    \boldsymbol\sigma
    =\lambda\,\operatorname{tr}(\boldsymbol\varepsilon)\mathbf I
     +2\mu\boldsymbol\varepsilon,
\]
где $\lambda$ и $\mu$ --- коэффициенты Ламе. Через модуль Юнга $E$ и
коэффициент Пуассона $\nu$
\[
    \mu=\frac{E}{2(1+\nu)},
    \qquad
    \lambda=\frac{E\nu}{(1+\nu)(1-2\nu)}.
\]
Подстановка закона Гука и кинематических соотношений в уравнение движения даёт
уравнения Ламе--Навье
\[
    \rho\ddot{\mathbf u}
    =\mu\Delta\mathbf u
     +(\lambda+\mu)\nabla(\nabla\!\cdot\!\mathbf u)
     +\rho\mathbf f.
\]
При статике левая часть равна нулю. Таким образом, в линейной теории упругости
неизвестным можно считать только поле перемещений $\mathbf u$, а деформации и
напряжения затем восстанавливаются по $\mathbf u$.

\subsection{Граничные и начальные условия}

Граница тела обычно разбивается на две части
$\partial\Omega=\Gamma_u\cup\Gamma_t$. На $\Gamma_u$ задаются кинематические
условия
\[
    \mathbf u=\overline{\mathbf u},
\]
а на $\Gamma_t$ --- силовые условия
\[
    \boldsymbol\sigma\mathbf n=\overline{\mathbf t}.
\]
Возможна и смешанная постановка, когда на разных частях границы задаются
различные компоненты перемещений и усилий. Для динамической задачи дополнительно
задаются начальные данные
\[
    \mathbf u(\mathbf x,0)=\mathbf u_0(\mathbf x),
    \qquad
    \dot{\mathbf u}(\mathbf x,0)=\mathbf v_0(\mathbf x).
\]

\subsection{Энергетическая и слабая постановки}

Для линейно-упругого тела с консервативной нагрузкой задача равновесия может
быть сформулирована как задача минимума полной потенциальной энергии
\[
    \Pi(\mathbf u)
    =\frac12\int_{\Omega}
       \boldsymbol\varepsilon(\mathbf u):\mathbf C:
       \boldsymbol\varepsilon(\mathbf u)\,d\Omega
     -\int_{\Omega}\rho\mathbf f\cdot\mathbf u\,d\Omega
     -\int_{\Gamma_t}\overline{\mathbf t}\cdot\mathbf u\,d\Gamma.
\]
Из условия $\delta\Pi=0$ получается слабая форма уравнений равновесия. Именно
эта постановка служит естественной основой методов Ритца, Галёркина и метода
конечных элементов.

\subsection{Неупругое деформирование}

Для реальных материалов одной линейной упругости недостаточно. В простейших
реологических моделях используются три базовых типа поведения:
\begin{itemize}
    \item \textbf{упругость} --- напряжение определяется текущей деформацией
    (элемент Гука);
    \item \textbf{вязкость} --- напряжение связано со скоростью деформации
    (элемент Ньютона);
    \item \textbf{пластичность} --- после достижения предела текучести возникают
    необратимые деформации (элемент Прандтля).
\end{itemize}
Комбинируя элементы, получают модели вязкоупругости. Например, в одномерном
случае модель Кельвина--Фойгта задаётся уравнением
\[
    \sigma=E\varepsilon+\eta\dot\varepsilon,
\]
а для модели Максвелла
\[
    \dot\varepsilon=\frac{\dot\sigma}{E}+\frac{\sigma}{\eta}.
\]
В пластичности обычно вводят разложение
$\boldsymbol\varepsilon=\boldsymbol\varepsilon^e+
\boldsymbol\varepsilon^p$, условие текучести и закон пластического течения.
Тем самым различия между моделями МДТТ в значительной мере определяются именно
выбором определяющих соотношений материала.

\subsection{Численное решение задач МДТТ}

После пространственной дискретизации, например методом конечных элементов,
непрерывная задача сводится к системе алгебраических или обыкновенных
дифференциальных уравнений. Для линейной статической задачи типична система
\[
    K\mathbf q=\mathbf F,
\]
где $K$ --- матрица жёсткости. В динамике возникает система
\[
    M\ddot{\mathbf q}+C\dot{\mathbf q}+K\mathbf q=\mathbf F(t).
\]
Для геометрически или физически нелинейных моделей решение обычно строится
пошагово по нагрузке или времени с итерационным уточнением на каждом шаге.

\subsection{Что важно сказать на экзамене}

Удобно представить модель МДТТ как четыре взаимосвязанных блока:
\emph{кинематика} $\mathbf u\mapsto\boldsymbol\varepsilon$,
\emph{баланс} $\operatorname{div}\boldsymbol\sigma+\rho\mathbf f
=\rho\ddot{\mathbf u}$, \emph{определяющий закон}
$\boldsymbol\sigma=\boldsymbol\sigma(\boldsymbol\varepsilon,\ldots)$ и
\emph{граничные/начальные условия}. Для линейной изотропной упругости нужно
записать тензор малых деформаций и закон Гука. Затем достаточно сказать, что
неупругие модели вводят вязкость и пластичность, а вариационная постановка
приводит к методу конечных элементов.

\newpage
